\subsection{2.2  The Performance Improvement Plan as a Structural Template}

The Performance Improvement Plan in organizational management is a formal document issued to an underperforming employee that specifies: the performance deficiencies observed, the root causes identified, the corrective actions required, the time window for improvement, the check-in schedule, and the criteria by which success or failure will be evaluated. The American Society for Human Resource Management recommends that a PIP should be used when there is a genuine commitment to help the employee improve, not merely as a precursor to termination [8]. Kirkpatrick [7, 8] emphasizes that the PIP should be developed collaboratively by the manager and employee, because remediation requires both parties' participation to succeed.

The PIP analogy to agent management is structurally productive but requires important qualification. A human employee can engage in self-directed improvement: they understand the feedback, adjust their behavior, and exercise agency. A deployed AI agent cannot. The APIP is therefore a management protocol imposed entirely by human operators---the agent is the subject, not the , of the plan. This distinction is theoretically important and has practical implications for where we locate the locus of accountability in the APIP lifecycle. In organizational terms, the APIP collapses the manager--employee dyad: the operator is both the manager who writes the plan and the agent of change who implements it.

